\section{Unconventional datasets}
Whilst classical procedures bank heavily upon either the assumption of Gaussianity in the exact sample domain, or employ Gaussianity results in the large sample domain, nonparametric procedures come up with distribution-invariant results with great generalization at the cost of reduced power and confidence. Still, both approaches struggle with phenomena such as bounded support, asymmetric and kurtotic concentration, non-polynomial tail decay, and multimodal tendencies. In many applications, forcing such datasets into Gaussian or exponentially-tailed frameworks may lead to inaccurate uncertainty quantification, unstable inference, and poor finite-sample approximations.

We aim to strike a middle ground of balance, where the shape-flexibility of our \sindist{} leads to more flexible data-distributional and test-statistic distributional models. This aids in the analysis of unconventional, specifically skewed or platykurtic, datasets. Such data-fitting exercise will constitute the foundational basis for later procedures like testing and inference. 

An example of fitting the \sindist{} to data was demonstrated in \autoref{sec:msmfit} and \autoref{sec:cdffit}.